\section{The whole ordinary binary tail is log-concave}
This section concerns the FORWARD ordinary polynomial, not its reversal
and not its positive derivative.  Let
\[
 F_h=F_{1^\alpha2^h},\qquad\alpha\ge1,\quad h\ge0,
 \qquad f_n=\operatorname{Fib}(n+1).
\]
All sequences are zero-extended.  The following assertions hold for
every initial padding length, without a count or rank cutoff.

\begin{theorem}[Ordinary binary full-tail induction]
The actual ordinary coefficients satisfy
\[
 F_h(n)=f_n\ (0\le n\le h+1),\quad
 0\le F_h(n)\le f_n,\quad F_h(h+2)\le f_{h+2}-1.       \tag{1}
\]
For every centre $n\ge h+1$,
\[
 \mathfrak m_h(n):=F_h(n)^2-F_h(n-1)F_h(n+1)\ge0.      \tag{2}
\]
Moreover $\mathfrak m_h(n)\ge1$ whenever additionally
$F_h(n)>F_h(n-1)$.  Statement (2) includes the ENTIRE tail
beyond the ordinary mode, through the terminal support and zero extension.
\end{theorem}
\begin{proof}
The exact actual recurrence is $F_{h+1}=1+(z+z^2)F_h$.
For $h=0$, the sequence is the indicator of $[0,\alpha]$;
(1)--(2) and the conditional strict assertion follow directly.
The recurrence proves (1) by induction.  In particular the first
possible defect of $F_{h+1}$ is the sum of the old first defect
and one exact Fibonacci value, so at least one unit is missing.

At the new boundary centre $n=h+2$, the three coefficients are
$f_{h+1}, f_{h+2}$ and at most $f_{h+3}-1$.
Cassini therefore gives
\[
 \mathfrak m_{h+1}(h+2)
       \ge f_{h+1}+(-1)^{h+2}\ge1.                      \tag{3}
\]
For $h=0$ this lower bound is two; for $h\ge1$ use
$f_{h+1}\ge2$.  Every later centre $n\ge h+3$ is a
two-term window of four old consecutive values $A,B,C,D$:
\[
 (B+C)^2-(A+B)(C+D)
   =(B^2-AC)+(C^2-BD)+(BC-AD).                          \tag{4}
\]
The two old log-concavity assertions apply at centres $n-2,n-1$.
If $B>0$, multiplying them by $C,A$ shows $BC\ge AD$.
If $B=0$, interval support forces $C=D=0$, so the new minor
is zero.  Thus every term of (4) is nonnegative.
If the new coefficient grows, $C>A$; hence either $B>A$ or
$C>B$.  The corresponding old minor is at least one by the
induction hypothesis, and so is (4).  This proves the full
unbounded induction, rather than only an induction before old modes.
\end{proof}

\begin{corollary}[Relative Fibonacci LC on the whole rising prefix]
If $h\ge1$ and $R$ is the FIRST ordinary mode of $F_h$, the
sequence $F_h(n)/f_n$ is log-concave through $R$.
\end{corollary}
\begin{proof}
At $n\le h+1$, the two earlier/current values have their exact
Fibonacci values and the next one is bounded above by its reference;
this proves the relative minor, including the first-defect boundary.
For $h+2\le n\le R$, the actual ordinary prefix strictly grows.
Here one uses the established interval support of positive ordinary
slopes starting at two; the initial equality $F_h(0)=F_h(1)$
causes no exception.  Thus (2) has integer minor at least one.
Put $A=F_h(n-1),B=F_h(n),C=F_h(n+1)$ and
$r=f_{n-1},s=f_n,t=f_{n+1}$.  Cassini and (1) give
\[
 B^2rt-ACs^2
   =rt(B^2-AC)-AC(s^2-rt)\ge rt-AC\ge0.               \tag{5}
\]
This is the desired relative LC.  No relative LC is asserted
after the FIRST ordinary mode: a long constant ordinary plateau
would make that stronger statement false.
\end{proof}

\section{A new unbounded nonbinary property-C family}
Let $B=1^\alpha2^h$ with $h\ge1$, and let $R$ be its first
ordinary forward mode.  Append a capacity $L\ge3$, making the
actual sorted arm $Y=BL$.

\begin{theorem}[Every capacity at or beyond the old first mode]
If $L\ge R$, then $Y$ has complete property C: at every
positive-slope comparison rank $N$ and every shift $s\ge3$,
\[
           G_Y(N-s)D_Y(N-1)\ge G_Y(N-s-1)D_Y(N).        \tag{6}
\]
In particular $L\ge\max(3,2h)$ is a sufficient capacity bound
independent of the initial padding.  Counts, padding, capacities,
comparison ranks and shifts are all unbounded.
\end{theorem}
\begin{proof}
First the reference $r_n=\operatorname{Fib}(n)$ for $n\ge2$
has complete C.  For any fixed comparison choose a previously proved
one-jump infinite-binary reference with capacity larger than every
read rank.  Its slopes agree with $r$ through that comparison,
so its complete-C theorem applies.  This is finite-prefix agreement,
not an analytic limiting argument.  Its normalization is $r_2=1$.

Consider the temporary slope prefix
\[
                      \widehat D_n=F_h(n-1),\quad n\ge2.
\]
Through rank $R+1$, the preceding relative-LC corollary supplies
normalized relative LC of $\widehat D_n/r_n$.
Its first two normalized values are one, since
$F_h(1)=1$ and $F_h(2)=2$.
The proved normalized-relative-C inheritance theorem therefore
gives every C comparison of $\widehat D$ at $N\le R+1$.
All temporary cumulatives use their required shared atom one at
zero and shared zero slope at one.

The exact actual append identity, for $n\ge2$, is
\[
 \Delta F_Y(n)=F_h(n-1)-F_h(n-L-1).                    \tag{7}
\]
If $L\ge R+1$, the actual slopes agree with $\widehat D$
through $R+1$.  With $M=\max F_h$, (7) is always at most
$M$, and equals $M$ first at $R+1$.  Thus the new FIRST
positive-slope mode is exactly $R+1$.
Every increasing comparison lies through that mode by the established
actual positive-slope unimodality.  Its cumulative reads are also
inside the common prefix, so the temporary C comparisons prove (6).
Decreasing or flat current slopes are immediate from $G_Y(k)\ge
G_Y(k-1)\ge0$.

The exact boundary $L=R$ requires retaining the outgoing unit.
For $n\le R$ the same prefix agreement holds.  At $n=R+1$,
the actual slope is $M-1$, rather than the temporary value $M$.
For $n<R+1$, the integer first-mode property gives
$F_h(n-1)\le M-1$.
For $n\ge R+1$, if the entering coefficient in (7) is positive,
interval support makes the outgoing coefficient at least one;
otherwise the signed slope is nonpositive.  Thus every positive
slope is at most $M-1$, attained at $R+1$, and the new first
slope mode is at most $R+1$.

The only additional potentially increasing current edge is
$N=R+1$.  Its preceding slope agrees with $\widehat D$;
all cumulative reads have rank at most $N-3=R-2$ and also agree.
Replacing the current slope $M$ by $M-1$ increases each C minor
by the nonnegative earlier cumulative $G_Y(N-s-1)$.
Thus that edge is paid as well.  This proves the sharp-boundary
version $L\ge R$, not only $L>R$.

Finally the first ordinary mode of a binary arm satisfies $R\le2h$.
At $h=0$ its first mode is zero.  Under a binary append every
edge beyond old first mode plus two is nonpositive, because
$\Delta F_{h+1}(n)=F_h(n-1)-F_h(n-3)$ there.
Hence $R_{h+1}\le R_h+2$, and induction gives the count bound.
\end{proof}

\begin{corollary}[All-depth coupled consequences]
For an arm in this one-jump family against an arbitrary opposite
arm, the existing one-arm-C bridge proves tail nonrecovery at
EVERY overlap depth when both ordinary factors rise and the current
tail is nonnegative.  The usual OLD head is retained for all-head S.
If both counts are at least two, the ENTIRE proper zero-head S
comparison holds without OLD A2 or a phase condition, by the complete
one-arm-C zero-head bridge and its already proved phase forcing.
\end{corollary}

\paragraph{Verification and remaining domain.}
The new Lean file proves the strict first-defect boundary, the
growth-sensitive two-term LC window lemma, the FULL unbounded
ordinary-family tail induction, the rising-prefix relative-LC
implication, count-to-mode induction, current-slope C improvement
and support/outgoing-unit budget.  The actual word recurrence,
rectangular base, first-mode/support identification and application
of the already proved relative-C/zero-head theorems are the explicit
written assembly above.  The earlier normalized-relative-C and
one-jump reference kernels are reused without recounting.
The separate actual-word QA is finite corroboration, not the proof.
No end-to-end Lean actual-composition theorem is claimed.

Capacities $3\le L<R$ are outside this theorem.  The complete
first-missing-four family and the first unresolved general C count
six are still open outside the recorded regions.  This is a NEW
unbounded nonbinary C family and coupled-comparison resource;
it does not prove full A2 for lighter binary arms, full original
coefficient UI for every one-jump arm, or upgrade the existing
original least-counterexample bounds.
